%% kapa: exam=1 round=24 year=2013 period=2 form=A subject=appraisal_law subject_ko=부동산관계법규 seq=35 number=75
%% source: Q-Net 감정평가사 기출문제 (cst00309) artlSeq=5073802 file=extracted/1008/5073802/2013년 제24회 감정평가사 1차시험 기출문제/2교시(A형).hwp
%% conversion: kordoc 4.24.51 -> tools/gen1.py (auto); review=proofread-sonnet; flags=none
\begin{question}{75}
측량ㆍ수로조사 및 지적에 관한 법령상 지목에 관한 설명으로 옳지 않은 것은?
\begin{choices}
  \item 사과ㆍ배 등 과수류를 집단적으로 재배하는 토지와 이에 접속된 저장고 등 부속시설물의 부지(주거용 건축물의 부지는 제외)는 “과수원”으로 한다.
  \item 축산업 및 낙농업을 하기 위하여 초지를 조성한 토지와 이에 접속된 부속시설물의 부지(주거용 건축물의 부지는 제외)는 “목장용지”로 한다.
  \item 천일제염 방식으로 하지 아니하고 동력으로 바닷물을 끌어들여 소금을 제조하는 공장시설물의 부지는 “염전”으로 한다.
  \item 「국토의 계획 및 이용에 관한 법률」 등 관계 법령에 따른 택지조성공사가 준공된 토지는 “대”로 한다.
  \item 제조업을 하고 있는 공장시설물의 부지와 같은 구역에 있는 의료시설 등 부속시설물의 부지는 “공장용지”로 한다.
\end{choices}
\end{question}
\answer{3}
